\chapter{پروسودی و ویژگی‌های زبرزنجیری گفتار}

در جلسات پیشین گفتار را از دید «چه گفته شد؟» بررسی کردیم: واج‌ها، فرمنت‌ها و \eng{MFCC}. این جلسه به پرسش مکمل می‌پردازد: «چگونه گفته شد؟». دو گوینده می‌توانند دقیقاً یک دنبالهٔ واجی را بگویند، اما یکی آن را خبری و دیگری پرسشی، با اطمینان یا با تردید، سریع یا آهسته ادا کند. این لایهٔ دوم را \term{پروسودی}{Prosody} می‌نامیم.

%==========================================================
\section{مرور کوتاه جلسهٔ قبل}
%==========================================================

\textbf{مراحل \eng{MFCC}} روی هر قاب کوتاه: تبدیل فوریه و طیف توان $\leftarrow$ بانک فیلتر مل $\leftarrow$ لگاریتم انرژی هر فیلتر $\leftarrow$ \eng{DCT} $\leftarrow$ نگه‌داشتن چند ضریب نخست (مثلاً ۱۲). فیلترهای مل در فرکانس‌های پایین باریک و در فرکانس‌های بالا پهن‌ترند. تعداد خروجی‌های بانک فیلتر برابر تعداد فیلترهاست (مثلاً ۴۰)، ولی تعداد ضرایب نگه‌داشته‌شده (مثلاً ۱۲) انتخاب جداگانهٔ طراح است؛ این دو الزاماً برابر نیستند.

\begin{correctionbox}{گام چهارم \eng{MFCC} تبدیل فوریهٔ دوباره نیست}
در توضیح شفاهی کلاس از «تبدیل مجدد» سخن رفت. در روش استاندارد، این گام \term{تبدیل کسینوسی گسسته}{DCT} است که روی لگاریتم انرژی‌های بانک فیلتر اعمال می‌شود و ضرایبی کم‌هم‌بسته‌تر می‌دهد.
\end{correctionbox}

\textbf{\eng{MFCC} و مدل منبع–فیلتر.} در مدل \term{منبع–فیلتر}{Source--Filter}، منبع (برای گفتار واک‌دار) ارتعاش تارهای صوتی با فرکانس پایهٔ \eng{$F_0$} است و فیلتر \emph{کل} مجرای صوتی از حنجره تا لب (نه فقط دهان، فک و دندان). \eng{MFCC} عمدتاً شکل طیفی و ویژگی‌های فیلتر را نگه می‌دارد و اطلاعات دقیق منبع، مثل \eng{$F_0$} و زیر و بمی، در آن تضعیف یا حذف می‌شود؛ زیرا فیلترهای مل ساختار هارمونیک را هموار می‌کنند و ضرایب نخست فقط پوش کند طیف را نگه می‌دارند. پس \eng{MFCC} برای تشخیص واج و طبقه‌بندی مناسب است (مثلاً طبقه‌بندی صدای سالم و بیمار، یا ورودی \eng{CNN})، اما برای مطالعهٔ آهنگ جمله کافی نیست؛ این انگیزهٔ بحث پروسودی است.

\textbf{مرور «آ» و «اَ».} تفاوت این دو واکه مربوط به \textbf{فیلتر} است، نه منبع: هر دو واک‌دارند و منبع یکسانی دارند، ولی موقعیت زبان فرق می‌کند و فرمنت‌ها جابه‌جا می‌شوند. هر دو زبان پایینی‌اند (\eng{$F_1$} نسبتاً بالا)، اما «آ» عقبی است و «اَ» جلویی، پس \eng{$F_2$} در «اَ» بالاتر است.

\begin{exambox}{نکتهٔ امتحانی}
هرچه زبان پایین‌تر باشد \eng{$F_1$} بالاتر است؛ هرچه زبان جلوتر باشد \eng{$F_2$} بالاتر است. باید بتوانید با \eng{$F_1$} و \eng{$F_2$} دربارهٔ ارتفاع و موقعیت افقی زبان در تولید واکه‌ها توضیح بدهید.
\end{exambox}

%==========================================================
\section{پروسودی چیست؟}
%==========================================================

\begin{definitionbox}{پروسودی (نوای گفتار)}
\term{پروسودی}{Prosody} مجموعهٔ ویژگی‌هایی از گفتار است که روی دنبالهٔ آواها و واژه‌ها «سوار» می‌شوند و به آن‌ها آهنگ، ریتم، تکیه و الگوی زمانی می‌دهند. پروسودی نشان می‌دهد \textbf{چگونه} گفته شده است، نه فقط \textbf{چه} گفته شده است.
\end{definitionbox}

مثال: «رفتی.» (خبری)، «رفتی؟» (پرسشی) و «رفتی!» (تعجب): واج‌ها یکسان، اما معنا و نیت گوینده متفاوت است. علائم نگارشی بخشی از این اطلاعات را در متن نشان می‌دهند، ولی تصویر کامل پروسودی در متن نیست و باید سیگنال صوتی را تحلیل کرد.

\textbf{قیاس با عروض.} استاد برای نشان دادن «لایه‌ای» بودن پروسودی به عروض اشاره کرد: شعر علاوه بر واژه و معنا روی یک الگوی وزنی (هجاهای کوتاه و بلند) سوار است، مثلاً «به نام خداوند جان و خرد» با الگوی «فعولن فعولن فعولن فعل». برای تحلیل وزن، معنای واژه‌ها اهمیتی ندارد؛ دو عبارت با معنای کاملاً متفاوت می‌توانند وزن یکسان داشته باشند. در گفتار عادی هم واج‌ها یک لایه‌اند و نحوهٔ بیان آن‌ها لایهٔ دیگر. (درس وارد تقطیع شعر نمی‌شود؛ این فقط یک قیاس است.)

\subsection{ویژگی‌های زبرزنجیری}
پروسودی را \term{ویژگی‌های زبرزنجیری}{Suprasegmental Features} هم می‌نامند: «فراتر از قطعه (واج)»؛ مثل چتری که روی قطعات گسترده است. دو پیامد:
\begin{itemize}
  \item تحلیل در سطح چند هجا، عبارت یا کل جمله انجام می‌شود، نه در پنجره‌های ۲۵ میلی‌ثانیه‌ای یا یک واج.
  \item آنچه اهمیت دارد افت و خیزها \emph{نسبت به بخش‌های قبل و بعد} است، نه مقدار مطلق در یک لحظه. از این رو مفهوم \term{کانتور}{Contour} (تغییرات یک ویژگی در طول زمان) محوری است.
\end{itemize}

\subsection{سه مؤلفهٔ اصلی}
«عصبانیت» یا «تردید» خودشان کمیت صوتی نیستند؛ باید ببینیم چه چیزهایی در سیگنال تغییر کرده که شنونده آن را چنین درک می‌کند. این کمیت‌ها در سه گروه جای می‌گیرند: فرکانس، شدت و زمان.

\begin{figure}[htbp]
\centering
\resizebox{\textwidth}{!}{%
\begin{tikzpicture}[rtl, font=\small,
  nd/.style={draw, rounded corners=6pt, align=center, text width=3.1cm, inner sep=5pt, line width=0.6pt},
  note/.style={align=center, text width=5.2cm, font=\footnotesize, text=black!70},
  arr/.style={-{Stealth[length=2.4mm]}, thick, mainblue},
  sk/.style={red!75!black, thick, -{Stealth[length=2mm]}}]
% centre
\node[draw, diamond, aspect=2.2, fill=softred, align=center, inner sep=2pt, line width=0.6pt] (sup) at (0,0) {\lr{suprasegmental}\\\lr{phonetic features}};
\node[note] (supn) at (0,2.9) {\lr{Supra}: فراتر؛ \lr{Segment}: قطعه (بخش). چتری که از بالا روی قطعات گسترده شده است؛ تحلیل در سطح کل جمله یا عبارت، نه فقط یک پنجره یا یک واج؛ افت و خیزها نسبت به بخش قبل و بعد.};
% duration
\node[nd, fill=yellow!30] (dur) at (-6.8,0) {\textbf{\lr{Duration}}\\زمان\\ریتم / مکث\\\lr{punctuation} (کاما)\\\lr{phrasing}\\رفع ابهام نحوی\\تعداد هجا در ثانیه\\کشیده‌شدن};
% intensity
\node[nd, fill=green!20] (int) at (6.3,0) {\textbf{\lr{intensity contour}}\\\lr{energy}\\شدت\\بلندی\\\lr{emphasis}};
\node[note, text width=5cm, font=\small] at (6.3,2.0) {\lr{$\text{Intensity}\propto\text{Magnitude}^{2}$}};
% pitch
\node[nd, fill=cyan!20] (pit) at (0,-3.4) {\textbf{\lr{Pitch Contour}}\\فرکانس\\زیر و بمی\\آهنگ کلام\\\lr{intonation}};
% pitch children
\node[nd, fill=cyan!20, text width=3.4cm] (dec) at (-3.1,-8.2) {کانتور پیچ جمله خبری: زیر و بمی صدا در آخر به سمت بم شدن می‌رود (سقوط فرکانس). نشان‌دهندهٔ قطعیت و پایان است.};
\node[nd, fill=cyan!20, text width=3.4cm] (que) at (0.9,-8.2) {کانتور پیچ جمله سؤالی: بعد از اکسنت، تارهای صوتی کشیده‌تر شده و فرکانس به شدت بالا (زیر) می‌رود («رفتی؟»).};
% intensity children
\node[nd, fill=green!20, text width=3.1cm] (str) at (4.4,-4.6) {\textbf{\lr{stress}}\\یک هجا در درون یک کلمه\\مثل: «مَردی» [اسم مصدر] در برابر «مردی» [یک مرد]};
\node[nd, fill=green!20, text width=3.4cm] (acc) at (8.4,-4.6) {\textbf{\lr{Accent/Focus}}\\یک کلمه در درون یک جمله\\کانتور اینتنسیتی: اکسنت = قله؛ سکوت یا گفتار بی‌تأکید = صاف};
\node[note, text width=4.2cm] (accn) at (8.4,-8.9) {کلمات در فرهنگ لغت و برای لهجهٔ خاص، استرس مشخصی دارند؛ اما این گوینده است که در موقعیت‌های مختلف تصمیم می‌گیرد به کدام کلمه در گفتمان (\lr{Discourse}) اکسنت بدهد.};
% arrows
\draw[arr] (sup.west) -- (dur.east);
\draw[arr] (sup.east) -- (int.west);
\draw[arr] (sup.south) -- (pit.north);
\draw[arr] (pit.south) -- ++(0,-0.4) -| (dec.north);
\draw[arr] (pit.south) -- ++(0,-0.4) -| (que.north);
\draw[arr] (int.south) -- ++(0,-0.4) -| (str.north);
\draw[arr] (int.south) -- ++(0,-0.4) -| (acc.north);
\draw[densely dashed, black!50] (acc.south) -- (accn.north);
% red sketches: fall / rise
\draw[sk] ([xshift=-0.9cm,yshift=-0.2cm]dec.south) -- ++(0.7,-0.8);
\node[red!75!black, font=\footnotesize] at ([xshift=0.8cm,yshift=-0.55cm]dec.south) {افت};
\draw[sk] ([xshift=-0.9cm,yshift=-0.2cm]que.south) to[out=-80,in=180] ++(0.4,-0.7) to[out=0,in=-110] ++(0.8,0.35);
\node[red!75!black, font=\footnotesize] at ([xshift=0.9cm,yshift=-0.55cm]que.south) {خیز};
\end{tikzpicture}}
\caption{نقشهٔ ذهنی ویژگی‌های زبرزنجیری: سه شاخهٔ زمان (\eng{Duration})، زیر و بمی (\eng{Pitch Contour}) و شدت (\eng{Intensity Contour}). شاخهٔ شدت به تکیهٔ واژگانی (\eng{stress}) و تأکید (\eng{Accent/Focus}) تقسیم شده است؛ شاخهٔ زیر و بمی به کانتور جملهٔ خبری و پرسشی. (بازسازی‌شده از تختهٔ کلاس)}
\end{figure}

\begin{table}[htbp]
\centering
\caption{سه مؤلفهٔ اصلی پروسودی، کمیت فیزیکی و ادراک متناظر. (جدول افزوده)}
\begin{tabular}{p{2.3cm}p{3.4cm}p{3.4cm}p{4.2cm}}
\toprule
مؤلفه & کمیت فیزیکی & ادراک شنونده & واحد/ابزار \\
\midrule
فرکانس & فرکانس پایه \eng{($F_0$)} & زیر و بمی (\eng{Pitch}) & هرتز، نیم‌پرده؛ کانتور \eng{$F_0$} \\
شدت & دامنه، انرژی & بلندی (\eng{Loudness}) & دسی‌بل؛ کانتور شدت \\
زمان & مدت آوا/هجا، مکث & سرعت، ریتم، کشش & میلی‌ثانیه، هجا در ثانیه \\
\bottomrule
\end{tabular}
\end{table}

%==========================================================
\section{مؤلفهٔ اول: شدت صدا}
%==========================================================

\subsection{شدت و دسی‌بل}
اصطلاحات \eng{Intensity}، \eng{Energy}، \eng{Loudness}، \eng{Amplitude} و \eng{Magnitude} از نظر علمی دقیقاً هم‌معنا نیستند، ولی در این جلسه در یک گروه قرار گرفته‌اند. رابطهٔ تخته: شدت با مربع دامنه متناسب است، $\text{\lr{Intensity}}\propto\text{\lr{Magnitude}}^{2}$. بلندی (\eng{Loudness}) کمیتی \textbf{ادراکی} است و با شدت فیزیکی یکی نیست. چون انسان نسبت‌های شدت را حس می‌کند و بازهٔ شدت‌های شنیدنی بسیار گسترده است، از مقیاس لگاریتمی \term{دسی‌بل}{dB} استفاده می‌شود.

\begin{definitionbox}{تراز شدت و تراز فشار صوت}
\begin{align*}
L_I &= 10\log_{10}\!\left(\frac{I}{I_0}\right)\ \text{dB}, \qquad I_0=10^{-12}\ \text{W/m}^2 \\
L_p &= \text{\lr{SPL}} = 20\log_{10}\!\left(\frac{p}{p_0}\right)\ \text{dB}, \qquad p_0=20\ \mu\text{Pa}
\end{align*}
چون $I\propto p^2$، ضریب لگاریتم برای دامنه ۲۰ و برای توان یا شدت ۱۰ است.
\end{definitionbox}

مثال‌ها: دو برابر شدن \textbf{دامنه} برابر $20\log_{10}2\approx+6\,\text{dB}$ است (نصف شدن: $-6\,\text{dB}$)؛ دو برابر شدن \textbf{توان} برابر $+3\,\text{dB}$؛ ده برابر شدن دامنه برابر $+20\,\text{dB}$. فشار $p=0{,}02\,\text{Pa}$ معادل $20\log_{10}(1000)=60\,\text{dB SPL}$ و در حد گفت‌وگوی معمولی است.

\coursefigure{0.85}{ch06-db-amplitude.png}{سمت چپ: یک تن خالص با دامنه‌های ۱، ۲ و ۱۰. سمت راست: رابطهٔ لگاریتمی دسی‌بل با نسبت دامنه؛ دو برابر شدن دامنه حدود $+6$ دسی‌بل و ده برابر شدن آن $+20$ دسی‌بل است. (شکل افزوده)}

\textbf{کانتور شدت:} برای هر قاب (مثلاً ۲۵ میلی‌ثانیه، گام ۱۰ میلی‌ثانیه) انرژی میانگین $E=\frac1N\sum_n x^2[n]$ را حساب و $10\log_{10}(E+\varepsilon)$ را در طول زمان رسم می‌کنیم.

\begin{correctionbox}{بلندی با شدت یکی نیست}
بلندی ادراک‌شده تابع شدت، فرکانس و مدت صوت است: دو صدا با شدت یکسان ولی فرکانس متفاوت بلندی یکسان ندارند و گوش در حدود ۲ تا ۵ کیلوهرتز حساس‌ترین است. قاعدهٔ سرانگشتی: حدود ۱۰ دسی‌بل افزایش، بلندی را تقریباً دو برابر می‌کند.
\end{correctionbox}

\subsection{تکیهٔ واژگانی و تأکید}

\begin{definitionbox}{\eng{Stress} در برابر \eng{Accent}}
\textbf{تکیهٔ واژگانی} (\eng{Lexical Stress}): برجستگی یک \textbf{هجا درون یک واژه}؛ بخشی از الگوی تلفظ واژه و وابسته به زبان و لهجه.

\textbf{تأکید} (\eng{Accent/Focus}): برجسته‌کردن یک \textbf{واژه درون جمله}؛ گوینده بر اساس قصد و بافت تصمیم می‌گیرد کدام واژه را برجسته کند. در کانتور شدت، تأکید به‌صورت یک \textbf{قله} دیده می‌شود و گفتار بی‌تأکید یا سکوت تقریباً صاف است.
\end{definitionbox}

مثال‌های کلاس: در «من گفتم که این کار را بکن» تأکید روی «من» است و در «کتاب را بردار» روی «کتاب».

\begin{warningbox}{«اکسنت» دو معنا دارد}
در این درس \eng{accent} یعنی «تأکید روی یک واژه در جمله»، نه «لهجه» (مثل لهجهٔ شیرازی)؛ و \eng{stress} یعنی «تکیه»، نه «استرس روانی».
\end{warningbox}

\textbf{رابطهٔ \eng{Accent} با \eng{Stress}} (نکته‌ای که استاد چند بار تکرار کرد): هنگام تأکید روی یک واژه، برجستگی روی همان هجایی قرار می‌گیرد که تکیهٔ واژگانی دارد؛ تأکید محل تکیه را عوض نمی‌کند و فقط برجستگی هجای تکیه‌دار را تقویت می‌کند. («نمی‌توان روی هجایی اکسنت گذاشت که استرس ندارد.») این ایده با \term{شبکهٔ متریک}{Metrical Grid} نمایش داده می‌شود؛ استاد گفت این مبحث بیشتر مربوط به تحلیل‌های زبان‌شناختی است و در درس کاربرد تخصصی ندارد.

\begin{correctionbox}{نشانگرهای تأکید فقط شدت نیستند}
کلاس تأکید را در شاخهٔ «شدت» جای داد و با قله در کانتور شدت شناخت؛ این درست ولی ناقص است. در پژوهش‌های آوایی، نشانگرهای تکیه و تأکید ترکیبی‌اند: تغییر \eng{$F_0$}، افزایش مدت هجا و کامل‌تر ادا شدن واکه، و سپس شدت (که در بسیاری از آزمایش‌ها ضعیف‌ترین نشانهٔ مستقل است). پس باید کانتور شدت، \eng{$F_0$} و مدت را هم‌زمان بررسی کرد.
\end{correctionbox}

\coursefigure{0.85}{ch06-synthetic-intensity.png}{کانتور شدت برای جملهٔ «کتاب را به علی دادی» در خوانش خنثی (خاکستری) و در حالتی که روی «علی» تأکید شده است (سبز): تأکید به‌صورت قله روی هجای تکیه‌دار آن واژه دیده می‌شود. داده‌ها مصنوعی و فقط برای توضیح مفهوم‌اند. (شکل افزوده)}

\begin{assignmentbox}{تمرین پیشنهادی با \eng{Praat}}
یک واژه یا جمله را یک بار عادی و یک بار با تأکید یا لحن عصبانی ضبط کنید و کانتور شدت و کانتور \eng{$F_0$} را مقایسه کنید. بررسی کنید که افزایش شدت و کشش در همان هجای تکیه‌دار رخ می‌دهد یا نه، و محل تکیه تغییر می‌کند یا نه.
\end{assignmentbox}

%==========================================================
\section{مؤلفهٔ دوم: فرکانس و زیر و بمی}
%==========================================================

\subsection{\eng{$F_0$}، \eng{Pitch} و آهنگ}
\begin{itemize}
  \item \term{فرکانس پایه}{\eng{$F_0$}}: کمیتی \textbf{فیزیکی}؛ تعداد دورهٔ ارتعاش تارهای صوتی در ثانیه (هرتز).
  \item \term{زیر و بمی}{Pitch}: کمیتی \textbf{ادراکی}؛ آنچه شنونده «زیر» یا «بم» می‌شنود. در عمل با \eng{$F_0$} یکی گرفته می‌شود.
  \item \term{کانتور زیر و بمی}{Pitch Contour}: تغییر \eng{$F_0$} در طول یک جمله یا عبارت.
  \item \term{آهنگ جمله}{Intonation}: الگوی منظم و معنادار این تغییرات که نقش زبانی دارد (مثلاً تمایز خبر و پرسش).
\end{itemize}

\begin{addedbox}{نیم‌پرده}
شنوایی زیر و بمی را تقریباً لگاریتمی حس می‌کند؛ پس برای مقایسه از \term{نیم‌پرده}{Semitone} استفاده می‌شود: $\Delta=12\log_2(f_2/f_1)$. مثلاً ۲۰۰ تا ۴۰۰ هرتز دقیقاً ۱۲ نیم‌پرده (یک اکتاو) است و ۲۰۰ تا ۲۵۰ هرتز حدود $3{,}9$ نیم‌پرده. همچنین \eng{$F_0$} فقط در قطعات واک‌دار تعریف می‌شود، پس کانتور در بی‌واک‌ها «پاره» است.
\end{addedbox}

\subsection{کانتور جملهٔ خبری و پرسشی}
\begin{itemize}
  \item \textbf{جملهٔ خبری}: زیر و بمی در پایان به سمت بم‌شدن می‌رود (سقوط فرکانس)؛ نشان‌دهندهٔ \textbf{قطعیت و پایان}.
  \item \textbf{جملهٔ پرسشی}: بعد از اکسنت، تارهای صوتی کشیده‌تر می‌شوند و فرکانس به شدت بالا می‌رود (مثال: «رفتی؟»).
\end{itemize}

\boardfigure{0.95}{ch06-board-metrical-grid.png}{کانتور \eng{$F_0$} نرمال‌شده (برحسب هرتز) در طول زمان نرمال‌شده برای چهار حالت \eng{Statement}، \eng{Question}، \eng{Happy} و \eng{Sad}. جملهٔ خبری در پایان افت می‌کند، پرسشی در پایان بالا می‌رود، حالت شاد در سطح بالاتر و تقریباً هموار است و حالت غمگین ابتدا بالا و سپس با افت شدید تمام می‌شود. دو منحنی رنگی با نوارهای اطمینان خاکستری دو گروه یا شرایط را نشان می‌دهند؛ تعریف آن‌ها در اسلاید ذکر نشده است.}

\begin{correctionbox}{«خبری = افت، پرسشی = خیز» گرایش است، نه قانون}
خیز پایانی در پرسش‌های بله/خیر رایج است، ولی پرسش‌های دارای واژهٔ پرسشی (چرا، کجا، چه) اغلب افتان‌اند. در فارسی گفتاری، پرسش بله/خیر اغلب فقط با آهنگ ساخته می‌شود («رفتی؟» در برابر «رفتی.»). شکل دقیق کانتور به گوینده، زبان و تأکید هم بستگی دارد. فارسی تونال نیست (برخلاف ماندارین که زیر و بمی در سطح واژه معنا می‌سازد).
\end{correctionbox}

\textbf{مثال «کتاب را به علی دادی؟»:} استاد این جمله را چند بار با لحن‌های گوناگون گفت: در حالت پرسشی \eng{$F_0$} در پایان بالا می‌رود و در حالت خبری پایین می‌آید. یک بار خیز پایانی به اندازهٔ انتظار روشن نبود؛ یعنی یک الگوی پروسودیک همیشه یک‌شکل تولید نمی‌شود و باید نمونه‌های زیادی را بررسی کرد. اگر روی «علی» تأکید شود، علاوه بر شدت، کانتور \eng{$F_0$} هم در آن ناحیه تغییر محسوسی دارد؛ پس بهتر است کانتور شدت و زیر و بمی را هم‌زمان بررسی کنیم.

\textbf{جمله‌های دستوری و احساسی:} در دستور قاطع ممکن است افت پایانی شدیدتر از جملهٔ خبری معمولی باشد؛ در درخواست مؤدبانه («بفرمایید، بشینید.») الگو متفاوت و لحن ملایم‌تر است؛ و شادی و غم کانتور و سطح متفاوتی دارند (شکل بالا). الگوی زیر و بمی به شیوهٔ بیان و قصد ارتباطی هم وابسته است، نه فقط کلمات.

%==========================================================
\section{مؤلفهٔ سوم: زمان، ریتم و مکث}
%==========================================================

\begin{table}[htbp]
\centering
\caption{ویژگی‌های زمانی گفتار (شاخهٔ \eng{Duration} در نقشهٔ کلاس).}
\begin{tabular}{p{3.3cm}p{9.5cm}}
\toprule
ویژگی & توضیح \\
\midrule
\eng{Duration} & مدت‌زمان بیان یک آوا، هجا، واژه یا عبارت \\
\eng{Speech Rate} & سرعت گفتار؛ ساده‌ترین معیار «تعداد هجا در ثانیه» است \\
\eng{Rhythm} & الگوی زمانی و ریتمیک گفتار \\
\eng{Pause} & مکث و سکوت \\
\eng{Lengthening} & کشیده‌شدن آواها \\
\eng{Phrasing} & تقسیم گفتار به گروه‌ها و عبارت‌های آهنگین \\
\bottomrule
\end{tabular}
\end{table}

\begin{addedbox}{سرعت گفتار در برابر سرعت تلفظ}
\term{سرعت گفتار}{Speech Rate} مکث‌ها را در زمان کل حساب می‌کند و \term{سرعت تلفظ}{Articulation Rate} فقط زمان گفتار واقعی را. مثلاً اگر ۱۹ هجا در $4{,}5$ ثانیه ادا شود که $0{,}5$ ثانیه آن مکث است: $19/4{,}5\approx 4{,}2$ هجا در ثانیه (گفتار) در برابر $19/4\approx 4{,}75$ (تلفظ).
\end{addedbox}

\textbf{مکث و علائم نگارشی.} مکث می‌تواند معنای نحوی جمله را تغییر دهد و مرز دو عبارت را روشن کند. ویرگول، دونقطه و نقطه نشانهٔ محل مکث یا بخش‌بندی‌اند؛ ولی همهٔ علائم زمانی نیستند (علامت سؤال به آهنگ مربوط است).

\textbf{رفع ابهام نحوی.} گاهی جمله در متن بیش از یک تفسیر دارد، اما در گفتار محل مکث‌ها، کشش آواها و گروه‌بندی زمانی ساختار مورد نظر را نشان می‌دهد (\term{رفع ابهام نحوی}{Syntactic Disambiguation}). مثلاً «مردان و زنان باهوش» (مثال افزوده، نه از کلاس): مکث پس از «مردان» یعنی «(مردان) و (زنان باهوش)» و مکث پس از «زنان» یعنی «(مردان و زنان) باهوش».

\coursefigure{0.85}{ch06-pause-phrasing.png}{دو گروه‌بندی ممکن برای «مردان و زنان باهوش» که فقط در جای مکث تفاوت دارند. مدت‌ها فرضی و برای توضیح‌اند. (شکل افزوده)}

%==========================================================
\section{پروسودی در تبدیل متن به گفتار}
%==========================================================

استاد چند نمونه خروجی از یک سامانهٔ فارسی تبدیل متن به گفتار پخش کرد. در نمونهٔ اول («مادر کتاب را روی میز اتاق گذاشت و پنجره را باز کرد.») مکث و گروه‌بندی طبیعی نبود: مدل واج‌ها را ساخته بود اما پروسودی مناسب جمله را بازسازی نکرده بود. در نمونهٔ دوم (جمله‌ای دربارهٔ شرکت \eng{Anthropic}) کسرهٔ اضافه نامناسب تولید می‌شد و سرعت گفتار نسبتاً بالا بود. (استاد از داده‌های آموزشی آن سامانه مطمئن نبود، پس نتیجهٔ قطعی دربارهٔ شیوهٔ آموزش آن نمی‌توان گرفت.)

کیفیت \eng{TTS} فقط به تلفظ درست واج‌ها وابسته نیست؛ سامانه باید محل مکث‌ها، طول هجاها، ریتم، الگوی تأکید، اتصال طبیعی واژه‌ها (مثل کسرهٔ اضافه) و آهنگ را هم درست تولید کند. دلیل دشواری: نگاشت متن به پروسودی «یک‌به‌چند» است و قصد گوینده و احساس در متن نیامده‌اند.

%==========================================================
\section{پژوهش: اطمینان و صداقت ادراک‌شده از روی پروسودی}
%==========================================================

\begin{assignmentbox}{تمرین طراحی آزمایش}
آزمایشی طراحی کنید که نشان دهد صداقت و اطمینان یک گوینده، در نحوهٔ بیان او و ویژگی‌های پروسودیک گفتارش، الگوهای مشترکی دارند.
\end{assignmentbox}

نکتهٔ مهم: باید میان (۱) اطمینان یا صداقتی که \textbf{شنونده} از صدا \emph{برداشت} می‌کند و (۲) اطمینان یا صداقت \textbf{واقعی} گوینده فرق گذاشت. آزمایش ادراکی کلاس موضوع اول را می‌سنجد.

\textbf{فرضیه‌های اولیهٔ کلاس} (قابل آزمون، نه قاعدهٔ قطعی):
\begin{itemize}
  \item \textbf{شدت}: شدت بیشتر یا تأکید واضح‌تر احساس تسلط می‌دهد؛ اما بلندی صدای افراد یکسان نیست، پس ابتدا باید ویژگی‌های معمول هر گوینده را به‌عنوان \term{خط پایه}{Baseline} در نظر گرفت و تغییرات همان فرد را بررسی کرد.
  \item \textbf{زیر و بمی}: افت \eng{$F_0$} در پایان عبارت قطعیت را منتقل می‌کند.
  \item \textbf{زمان و مکث}: مکث‌های متعدد در میانهٔ جمله تردید را می‌رساند؛ اما مکث در ابتدا یا انتها ممکن است متانت را برساند. پس \textbf{محل} مکث هم مهم است.
\end{itemize}

\textbf{طرح آزمایش پیشنهادی:}
\begin{center}
\begin{tikzpicture}[rtl, node distance=0.7cm and 0.9cm, >=Stealth,
  box/.style={draw, rounded corners, align=center, minimum width=2.7cm, minimum height=1.3cm, font=\small, fill=blue!6}]
  \node[box] (a) {\RL{۱. جمع‌آوری گفتار}\\\RL{جملهٔ یکسان با}\\\RL{شیوه‌های بیان متفاوت}};
  \node[box, left=of a] (b) {\RL{۲. ارزیابی شنونده}\\\RL{امتیاز اطمینان و}\\\RL{صداقت ادراک‌شده}};
  \node[box, left=of b] (c) {\RL{۳. استخراج ویژگی}\\\lr{Pitch, Intensity,}\\\lr{Duration}};
  \node[box, left=of c] (d) {\RL{۴. تحلیل آماری}\\\RL{هم‌بستگی ویژگی‌ها}\\\RL{با امتیاز شنوندگان}};
  \draw[->] (a) -- (b); \draw[->] (b) -- (c); \draw[->] (c) -- (d);
\end{tikzpicture}
\end{center}

\begin{definitionbox}{هم‌بستگی (\eng{Correlation})}
\textbf{هم‌بستگی مثبت}: با افزایش یک متغیر، دیگری نیز تمایل به افزایش دارد (مثلاً شدت بیشتر با امتیاز اطمینان بیشتر). \textbf{هم‌بستگی منفی}: دیگری تمایل به کاهش دارد (مثلاً مدت طولانی‌تر با اطمینان کمتر). نزدیک صفر یعنی رابطهٔ خطی ضعیف است. ضریب پیرسون بین $-1$ و $+1$ است و هم‌بستگی به‌تنهایی علیت را ثابت نمی‌کند.
\end{definitionbox}

\subsection{مقالهٔ بررسی‌شده}
مقاله با عنوان
\textit{\lr{Listeners' perceptions of the certainty and honesty of a speaker are associated with a common prosodic signature}}
(\eng{Goupil et al., Nature Communications, 2021}) است. در مطالعهٔ اول طبق گزارش‌ها، شنوندگان بین دو نسخه از یک شبه‌واژه که \eng{$F_0$}، بلندی و مدت آن‌ها تصادفی تغییر کرده بود انتخاب می‌کردند کدام «مطمئن‌تر» یا «صادق‌تر» است، و با \term{هم‌بستگی معکوس}{Reverse Correlation} «هستهٔ پروسودیک» هر قضاوت به دست می‌آمد. (آنچه روی تخته آمد خلاصهٔ آموزشی است؛ جزئیات را از متن مقاله بازبینی کنید.)

\textbf{خواندن نمودارها.} ردیف بالا مربوط به \textbf{اطمینان} (آبی) و ردیف پایین \textbf{صداقت} (سبز) است؛ ستون‌ها \eng{pitch}، \eng{loudness} و \eng{duration}. محور افقی «قطعه» است: بخش‌های زمانی یک محرک کوتاه (۱۲ قطعه برای \eng{$F_0$} و بلندی، ۵ قطعه برای مدت)، نه بخش‌های یک جملهٔ بلند. محور عمودی دامنهٔ هسته؛ خط خطچین صفر یعنی «بدون اثر».

\boardfigure{0.95}{ch06-paper-kernels.png}{هسته‌های پروسودیک اطمینان (ردیف بالا) و صداقت (ردیف پایین) برای \eng{pitch}، \eng{loudness} و \eng{duration} در برابر قطعهٔ زمانی. مقدار مثبت یعنی ویژگی در آن قطعه با انتخاب «مطمئن‌تر/صادق‌تر» همراه بوده؛ منفی یعنی عکس؛ نزدیک صفر یعنی بی‌اثر. (مقاله: \eng{Goupil et al., 2021})}

\begin{table}[htbp]
\centering
\caption{خلاصهٔ الگوهای نمودار.}
\begin{tabular}{p{2.1cm}p{6.1cm}p{5.2cm}}
\toprule
ویژگی & الگوی دیده‌شده & برداشت \\
\midrule
\eng{Pitch} & در ابتدا نزدیک صفر، در نیمهٔ دوم منفی & افت آهنگ در بخش پایانی با اطمینان و صداقت بیشتر همراه است \\
\eng{Loudness} & مثبت در چند قطعه (ابتدایی/میانی) & شدت بیشتر در این بخش‌ها با اطمینان یا صداقت بیشتر همراه است \\
\eng{Duration} & منفی در قطعه‌های میانی، نزدیک صفر در دو سر & طولانی‌تر (کندتر) بودن با اطمینان و صداقت کمتر همراه است \\
\bottomrule
\end{tabular}
\end{table}

نکتهٔ اصلی: اثر هر ویژگی به \textbf{جایگاه زمانی} آن بستگی دارد (مثلاً شدت در ابتدا بی‌اثر و در میانه مؤثر است)، پس میانگین کل جمله کافی نیست.

\begin{correctionbox}{دقت در تفسیر نمودار}
(۱) نمودار مدت به سرعت گفتار مربوط است، نه به «مکث» به معنای سکوت؛ در مقاله سرعت بیشتر (مدت کوتاه‌تر) با اطمینان بیشتر همراه است. (۲) «ابتدا، میانه و انتهای جمله» در کلاس را باید «ابتدا، میانه و انتهای محرک (واژه)» خواند. (۳) ایدهٔ مقاله «شکل» کانتور است؛ «فرکانس پایین‌تر» یعنی افت آهنگ در بخش پایانی، نه صدای بم‌تر به‌طور کلی. (۴) نتیجه دربارهٔ \textbf{برداشت شنونده} است، نه صداقت واقعی، و همبستگی است نه علیت.
\end{correctionbox}

%==========================================================
\section{کاربردهای عملی تحلیل پروسودی}
%==========================================================

یک نمونه که استاد اشاره کرد ارزیابی نحوهٔ بیان افراد در \textbf{مصاحبهٔ شغلی} است. چنین سامانه‌ای می‌تواند ترکیب کند: (۱) \textbf{محتوا} با \eng{ASR}، (۲) \textbf{نحوهٔ بیان} (شدت، فرکانس، سرعت گفتار) و (۳) \textbf{تحلیل چندوجهی} (متن، صدا و گاهی تصویر) برای ارائهٔ امتیاز یا برآورد.

\begin{warningbox}{احتیاط در استفاده}
چنین امتیازهایی لزوماً بیانگر شخصیت، تخصص یا صداقت واقعی فرد نیستند و اعتبار آن‌ها نیازمند ارزیابی مستقل و کنترل تفاوت‌های فردی، زبانی و موقعیتی است.
\end{warningbox}

%==========================================================
\section{جمع‌بندی}
%==========================================================

پروسودی لایه‌ای از گفتار است که «چگونه گفتن» را توصیف می‌کند و تمام آن در متن نیست. سه مؤلفهٔ آن زیر و بمی، شدت و زمان‌اند و باید به‌صورت کانتور روی بازه‌های بلند بررسی شوند؛ تحلیل طیفی به‌تنهایی کافی نیست.

\begin{exambox}{نکات مهم برای مرور و امتحان}
\begin{enumerate}
  \item \eng{MFCC} شکل طیفی و فیلتر مجرا را نشان می‌دهد و اطلاعات \eng{$F_0$} را عملاً حذف می‌کند. تعداد خروجی‌های بانک فیلتر برابر تعداد فیلترهاست و ضرایب نهایی (مثلاً ۱۲) می‌تواند کمتر باشد؛ گام پس از لگاریتم \eng{DCT} است.
  \item اختلاف «آ» و «اَ» مربوط به فیلتر است: زبان پایین‌تر $\Rightarrow$ \eng{$F_1$} بالاتر؛ زبان جلوتر $\Rightarrow$ \eng{$F_2$} بالاتر.
  \item پروسودی (زبرزنجیری) ریتم، آهنگ و تأکید را توصیف می‌کند؛ سه مؤلفه: \eng{Pitch/Frequency}، \eng{Intensity}، \eng{Time/Duration}؛ باید یک عبارت و کانتور آن را بررسی کرد.
  \item \eng{Stress}: تکیه درون واژه؛ \eng{Accent}: برجسته‌سازی واژه در جمله؛ تأکید روی همان هجای تکیه‌دار می‌نشیند.
  \item کانتور شدت ناحیه‌های تأکیدی را نشان می‌دهد، اما \eng{$F_0$} و مدت هم در تأکید نقش دارند. شدت با مربع دامنه متناسب است؛ $L_p=20\log_{10}(p/p_0)$ و $L_I=10\log_{10}(I/I_0)$.
  \item کانتور پرسشی اغلب با خیز و خبری با افت پایانی همراه است (گرایش، نه قانون).
  \item ویژگی‌های زمانی (سرعت گفتار، مدت هجا، ریتم، مکث، کشش) در \eng{Phrasing} و رفع ابهام نحوی نقش دارند.
  \item آزمایش اطمینان/صداقت: جمع‌آوری گفتار، ارزیابی شنونده، استخراج ویژگی، تحلیل هم‌بستگی؛ جایگاه زمانی ویژگی مهم است و باید بتوانید هم‌بستگی مثبت/منفی/نزدیک صفر را در نمودار تفسیر کنید.
\end{enumerate}
\end{exambox}
